\section{Reinforcement Learning Case of Study}\label{sec: rl case of study}

This section presents a case of study using the forward kinematics estimation in a \gls{rl} loop, which consists of two steps.
%
In the fist step, we train $\nn$ and a policy $\policy:\realSet{\nDims\times\nJoints}\rightarrow\realSet{\nJoints}$, which given the Cartesian position of the robot joints, learns the torques for each joint for moving the end-effector towards the target position, using the pre-trained $\nn$ model learned in Section \ref{sec: results} for estimating $\vector{\pos} = \nn(\vector{\rotAngle{}})$.
%
In the second step, we create a different robot with the same number of joints, by varying its links' dimension according to a normal distribution $\normalDist{\linksLenght}{\var=0.01}$, where $\linksLenght$ is the dimensions of the robot used in the first step, and retrain both $\nn$ and $\policy$ on the new robot.
%
As such, $\policy$ encompasses learning the robot's dynamics, and control, which is a considerably hard task for high-dimensions.

This study case mimics a straight-forward transfer-learning case, in which first a model is firstly trained in simulation, and then refined in a real robot, which is slightly different from the simulated one.
%
These differences may come from imprecise measures, or from an unknown model, as in the case of budget robot manipulators \cite{zajc2022low}. 
%
Furthermore, more elaborated transfer-learning approaches \cite{stuber2018feature,bocsi2013alignment} might also benefit from the proposed network, but its evaluation falls out-or-the-scope with this paper.

The results in this section are limited to the 3D and a robot with $7$ joints, and the second part of the learning is also done in simulation. 
%
%This limitation is due to the extensive long time necessary to tune, and learn $\policy$, and then to adapt both $\nn$ and $\policy$ to the new environment, since in the \gls{rl} loop, the reward is sparse and the robot performs a random-walk sampling.

For $\policy$ we use a \sac from \sb \cite{stablebaselines3}, and $\nn$ is introduced into the \sac as such its outputs are the inputs of \sac's features extraction step. 
Figure \ref{fig: learn and adapt policy} presents the rewards obtained in both steps, where we can see...

\begin{figure}
	\includegraphics[width=\linewidth]{example-image-a}
	\caption{Reward for learning $\policy$ using the pre-trained $\nn$, and for adapting $\policy$ to the modified robot using.}
	\label{fig: learn and adapt policy}
\end{figure}
